% chktex-file 36
\section{Chức năng Trace}

\subsection{Mục tiêu}
Thêm system call \texttt{trace(mask)} để theo dõi các system call của một tiến trình. Mỗi bit \texttt{n} trong \texttt{mask} tương ứng với system call có số hiệu \texttt{n} (ví dụ \texttt{1 << SYS\_fork}). Khi một system call nằm trong \texttt{mask} trả về, kernel in một dòng gồm pid của tiến trình gọi, tên system call và giá trị trả về. Việc theo dõi áp dụng cho tiến trình gọi \texttt{trace} và mọi tiến trình con do nó \texttt{fork} sau đó, không ảnh hưởng tiến trình khác.

\subsection{Tổ chức triển khai}
\begin{itemize}
  \item \texttt{kernel/syscall.h}: thêm \texttt{SYS\_trace} (số 22).
  \item \texttt{user/usys.pl}, \texttt{user/user.h}: thêm stub và nguyên mẫu \texttt{int trace(int)} để chương trình người dùng gọi được.
  \item \texttt{kernel/proc.h}: thêm trường \texttt{trace\_mask} vào \texttt{struct proc}.
  \item \texttt{kernel/sysproc.c}: cài đặt \texttt{sys\_trace()}, lấy đối số bằng \texttt{argint()} và lưu vào \texttt{trace\_mask} của tiến trình hiện tại.
  \item \texttt{kernel/proc.c}: \texttt{freeproc()} đặt lại \texttt{trace\_mask = 0}; \texttt{fork()} sao chép \texttt{trace\_mask} từ cha sang con.
  \item \texttt{kernel/syscall.c}: đăng ký \texttt{sys\_trace} vào mảng \texttt{syscalls}, thêm mảng tên \texttt{syscall\_names} và in kết quả trong \texttt{syscall()}.
  \item \texttt{user/trace.c}: chương trình người dùng \texttt{trace mask command [args...]}; thêm \texttt{\$U/\_trace} vào \texttt{UPROGS} trong \texttt{Makefile}.
\end{itemize}

\subsection{Luồng xử lý}
Lệnh \texttt{trace 32 grep hello README} chạy chương trình \texttt{trace}. Chương trình gọi \texttt{trace(32)} để kernel ghi mask vào \texttt{struct proc} của chính nó, rồi gọi \texttt{exec()} để thay thế bằng \texttt{grep}. Vì \texttt{exec()} giữ nguyên \texttt{struct proc} nên mask vẫn còn hiệu lực với \texttt{grep}.

Mỗi lần có system call, \texttt{syscall()} gọi hàm xử lý rồi kiểm tra \texttt{p->trace\_mask \& (1 << num)}. Nếu bit được đặt, kernel in \texttt{<pid>: syscall <tên> -> <giá trị trả về>}. Việc in được thực hiện \emph{sau} khi system call chạy xong để có giá trị trả về. Với \texttt{exit} (không bao giờ trả về) sẽ không có dòng nào được in.

\subsection{Khó khăn và giới hạn}
\begin{itemize}
  \item Phải đặt lại \texttt{trace\_mask} trong \texttt{freeproc()}, nếu không một tiến trình mới có thể dùng lại slot cũ và vô tình thừa hưởng mask.
  \item Mask là kiểu \texttt{int} nên chỉ theo dõi được system call có số hiệu nhỏ hơn 32 (hiện có 22 system call).
  \item Không in đối số của system call, theo đúng yêu cầu đề bài.
\end{itemize}

\subsection{Kiểm thử}
Ghi nhận kết quả chạy thực tế trên xv6 (pid có thể khác giữa các lần chạy; các giá trị trả về của \texttt{read} phụ thuộc nội dung file \texttt{README}):

\begin{center}
\begin{minipage}{0.86\textwidth}
\begin{lstlisting}[numbers=none,xleftmargin=0pt,framexleftmargin=0pt]
$ trace 32 grep hello README
3: syscall read -> 438
3: syscall read -> 0
$ trace 2147483647 grep hello README
4: syscall trace -> 0
4: syscall exec -> 3
4: syscall open -> 3
4: syscall read -> 1023
4: syscall read -> 991
4: syscall read -> 438
4: syscall read -> 0
4: syscall close -> 0
$ grep hello README
$ trace 2 usertests forkforkfork
usertests starting
6: syscall fork -> 7
test forkforkfork: 6: syscall fork -> 8
8: syscall fork -> 9
9: syscall fork -> 10
...
\end{lstlisting}
\end{minipage}
\end{center}

\noindent Ba kết quả đầu khớp đề: chỉ theo dõi \texttt{read}; theo dõi tất cả system call; và không có đầu ra khi chạy \texttt{grep} trực tiếp. Ở ví dụ thứ tư, các system call \texttt{fork} của mọi tiến trình con cháu đều được in, chứng tỏ mask được kế thừa qua \texttt{fork}.
